\begin{mistake}{2026-04-09}{41}

\begin{problemcontent}
计算 \( \displaystyle\int_0^{\frac{\sqrt{2}}{2}R} e^{-y^2}\,\mathrm{d}y \int_0^y e^{-x^2}\,\mathrm{d}x + \int_{\frac{\sqrt{2}}{2}R}^{R} e^{-y^2}\,\mathrm{d}y \int_0^{\sqrt{R^2 - y^2}} e^{-x^2}\,\mathrm{d}x \)。
\end{problemcontent}

\begin{solutioncontent}
将被积函数统一写为 \( e^{-(x^2+y^2)} \)，分析两个累次积分对应的区域：
区域 \( D_1 \)：\( 0 \leqslant y \leqslant \frac{\sqrt{2}}{2}R \)，\( 0 \leqslant x \leqslant y \)。
区域 \( D_2 \)：\( \frac{\sqrt{2}}{2}R \leqslant y \leqslant R \)，\( 0 \leqslant x \leqslant \sqrt{R^2 - y^2} \)。

画图可知，\( D_1 \) 的边界 \( x=y \) 与 \( D_2 \) 的圆弧边界 \( x^2+y^2=R^2 \) 在第一象限的交点纵坐标恰好为 \( y = \frac{\sqrt{2}}{2}R \)。
因此，这两个区域无缝拼合成了一个位于第一象限、介于 \( y \) 轴（\( x=0 \)）和射线 \( y=x \) 之间的圆心角为 \( \pi/4 \) 的扇形区域 \( D \)。

转入极坐标系：
极角范围：\( \frac{\pi}{4} \leqslant \theta \leqslant \frac{\pi}{2} \)
极径范围：\( 0 \leqslant r \leqslant R \)
\[
\begin{aligned}
原式 &= \iint_D e^{-(x^2+y^2)} \,\mathrm{d}x\mathrm{d}y \\
&= \int_{\pi/4}^{\pi/2} \mathrm{d}\theta \int_0^R e^{-r^2} r\,\mathrm{d}r \\
&= \left( \frac{\pi}{2} - \frac{\pi}{4} \right) \left[ -\frac{1}{2}e^{-r^2} \right]_0^R \\
&= \frac{\pi}{4} \cdot \frac{1-e^{-R^2}}{2} = \frac{\pi(1 - e^{-R^2})}{8}
\end{aligned}
\]
\end{solutioncontent}

\mistakeknowledge{
\begin{itemize}
 \item \textbf{核心方法}：利用二重积分对区域的可加性逆向操作，合并具有相同被积函数的复杂累次积分，再转极坐标求解。
 \item \textbf{易错点}：直角坐标化极坐标时，容易遗漏面积微元中的雅可比行列式即额外的 \( r \)，导致 \( e^{-r^2} \) 无法积分。
 \item \textbf{关联知识}：凑微分技巧 \( \int e^{-r^2} r\,\mathrm{d}r = -\frac{1}{2} \int e^{-r^2} \,\mathrm{d}(-r^2) = -\frac{1}{2} e^{-r^2} \)。
\end{itemize}}

\end{mistake}
